\subsection{Signed coherent preparation of a Taylor remainder}
\label{sec:coherent-tilt-remainder}

I retain the classical probability model throughout. Quantum states in
this construction represent computational sampling distributions and do
not replace the financial risk-neutral measure. The following is a
project-specific derivation of an importance preparation; a literature
comparison is required before claiming an original quantum method.

\begin{proposition}[Continuous coherent tilt identity]
Let $Y\sim N(0,I_p)$, $w_i>0$, $\sum_iw_i=1$, and let
\[
 T(y)=\sum_{i=1}^n w_i e^{d_i+v_i\cdot y},\qquad U(y)=T(y)-1,
 \qquad\sum_iw_i d_i=0,\quad\sum_iw_i v_i=0.
\]
For an integer $s\geq1$, set $M_{2s}=\mathbb E U(Y)^{2s}$ and
$\psi_0(y)=(2\pi)^{-p/4}e^{-\|y\|^2/4}$. Assume $M_{2s}>0$.
For $j=0,\ldots,s$ and multi-indices $\alpha\in\mathbb N_0^n$
with $|\alpha|=j$, define
\[
 c_{j,\alpha}=(-1)^{s-j}\binom sj\frac{j!}{\alpha!}w^\alpha,
 \quad d_\alpha=\sum_i\alpha_i d_i,
 \quad v_\alpha=\sum_i\alpha_i v_i,
 \quad \beta_{j,\alpha}=c_{j,\alpha}e^{d_\alpha+\|v_\alpha\|^2}.
\]
Let $\psi_m(y)=(2\pi)^{-p/4}e^{-\|y-m\|^2/4}$ and
$L=\sum_{j,\alpha}|\beta_{j,\alpha}|$. Then
\[
 U^s\psi_0=\sum_{j,\alpha}\beta_{j,\alpha}\psi_{2v_\alpha},
 \qquad
 \psi_s=\frac{U^s\psi_0}{\sqrt{M_{2s}}},
 \qquad
 \eta=\frac{M_{2s}}{L^2}\in(0,1].
\]
There are $\binom{n+s}{s}$ labelled components before any exact merging.
\end{proposition}

\begin{proof}
Weighted Jensen gives $T\geq\exp(\sum_iw_i(d_i+v_i\cdot y))=1$,
so $U\geq0$. The binomial and multinomial formulae give
\[
 U^s=\sum_{j=0}^s(-1)^{s-j}\binom sj
       \sum_{|\alpha|=j}\frac{j!}{\alpha!}w^\alpha
       e^{d_\alpha+v_\alpha\cdot y}.
\]
Completing the square gives
$e^{v\cdot y}\psi_0(y)=e^{\|v\|^2}\psi_{2v}(y)$.
Substitution proves the coherent identity. Every component is normalised,
so the triangle inequality gives
$\sqrt{M_{2s}}=\|U^s\psi_0\|\leq L$. The density of $\psi_s$ is
$q_s(y)=U(y)^{2s}\phi_p(y)/M_{2s}$. Gaussian exponential moments
justify its finite normalisation for every fixed $s$. The number of
multi-indices at degree $j$ is $\binom{n+j-1}{j}$; summing these
counts gives $\binom{n+s}{s}$.
\end{proof}

\begin{proposition}[Finite-register construction]
Let the finite base distribution factor as
$p(y)=\prod_{r=1}^p p_r(y_r)$ on a finite product grid, with all
$p_r(y_r)>0$. Define
\[
 Z(v)=\mathbb E_p e^{2v\cdot Y}
      =\prod_{r=1}^p\sum_{y_r}p_r(y_r)e^{2v_r y_r},
 \qquad
 |\varphi_v\rangle=\sum_y
       \frac{\sqrt{p(y)}e^{v\cdot y}}{\sqrt{Z(v)}}|y\rangle,
\]
and $\beta^{\mathrm{fin}}_{j,\alpha}
=c_{j,\alpha}e^{d_\alpha}\sqrt{Z(v_\alpha)}$.
Let $P_\ell$ be an explicitly emitted unitary preparing
$|\varphi_{v_\ell}\rangle$ from zero, let
$L_{\mathrm{fin}}=\sum_\ell|\beta^{\mathrm{fin}}_\ell|$, and let
$P_\chi$ prepare
$|\chi\rangle=\sum_\ell\sqrt{|\beta^{\mathrm{fin}}_\ell|/
L_{\mathrm{fin}}}|\ell\rangle$.
Define
\[
 \mathrm{SELECT}=\sum_\ell |\ell\rangle\langle\ell|
          \otimes\operatorname{sign}(\beta^{\mathrm{fin}}_\ell)P_\ell,
 \qquad
 B=(P_\chi^\dagger\otimes I)\mathrm{SELECT}(P_\chi\otimes I).
\]
Padding indices have zero weight and any stated unitary extension.
Then the zero-selector block of $B|0\rangle$ is
\[
 \frac1{L_{\mathrm{fin}}}\sum_y\sqrt{p(y)}U(y)^s|y\rangle,
\]
and its probability is
$\eta_{\mathrm{fin}}=M_{2s}^{\mathrm{fin}}/L_{\mathrm{fin}}^2$.
\end{proposition}

\begin{proof}
The first preparation supplies amplitudes
$\sqrt{|\beta_\ell|/L}$. The selector-controlled sign and component
unitary produce that amplitude times
$\operatorname{sign}(\beta_\ell)|\varphi_{v_\ell}\rangle$.
Projecting the unprepared selector onto zero multiplies by
$\langle0|P_\chi^\dagger|\ell\rangle
=\sqrt{|\beta_\ell|/L}$. The resulting vector is
$L^{-1}\sum_\ell\beta_\ell|\varphi_{v_\ell}\rangle$.
Substitution of the finite tilt normalisation cancels $\sqrt Z$ and
leaves the same binomial expansion of $U^s$. Squaring its norm proves
the probability formula. The component signs are conditional phases;
discarding them would prepare a different vector and invalidate the
identity.
\end{proof}

\begin{proposition}[Financial correction identity]
Let $g_y(t)$ be the monetary conditioned payoff, with $T(y)=1+U(y)$.
Let the order $k=2s-1$ Taylor polynomial be $J_k(y)$ and its remainder
$R(y)=g_y(T(y))-J_k(y)$. Suppose the checked derivative certificate gives
\[
 |R(y)|\leq b_s U(y)^{2s},\qquad
 b_s=\frac{H D_{2s,\lambda}}{(2s)!}.
\]
Define $v(y)=R(y)/(b_sU(y)^{2s})$ when $U(y)>0$ and $v(y)=0$
when $U(y)=0$. Then $|v|\leq1$ and
\[
 \mathbb E_p g_Y(T(Y))
 =\mathbb E_pJ_k(Y)+b_sM_{2s}^{\mathrm{fin}}\mathbb E_{q_s}v(Y).
\]
For the continuous identity, replace the finite measures and moments
by their continuous counterparts.
\end{proposition}

\begin{proof}
The remainder bound implies $R=0$ at $U=0$. On the other points,
$q_s(y)v(y)=p(y)R(y)/(b_sM_{2s})$. Summing proves the identity.
The continuous version follows by integration; the same remainder bound
ensures integrability. A flag of probability $(1+v(y))/2$ estimates the
signed expectation, not the expectation of a squared amplitude.
The classical polynomial term and its arithmetic errors remain part
of the total price and error budget.

The apparent singular quotient is avoidable. Set $m=2s$. The integral
Taylor remainder gives
\[
 R(y)=\frac{U(y)^m}{(m-1)!}\int_0^1(1-t)^{m-1}
                  g_y^{(m)}(1+tU(y))\,dt.
\]
Consequently a bounded extension of the importance observable is
\[
 v(y)=\frac{m}{H D_{m,\lambda}}\int_0^1(1-t)^{m-1}
                  g_y^{(m)}(1+tU(y))\,dt.
\]
The factor $m(1-t)^{m-1}$ integrates to one. Since $1+tU\geq1$,
the derivative certificate bounds the integrand divided by $H D_m$
in absolute value by one. This is a Beta$(1,m)$ expectation and remains
bounded at $U=0$, where its continuous value is
$g_y^{(m)}(1)/(H D_m)$. The importance measure gives that event zero
weight, so either bounded definition there yields the same identity.
An implementation must charge the Beta preparation or finite quadrature,
its bias, and the derivative oracle. This identity eliminates a
conditioning problem but does not establish an inexpensive circuit.
\end{proof}

An integer number of ordinary amplitude-amplification iterations does
not in general prepare $q_s$ with probability one. If
$\theta=\arcsin\sqrt\eta$ and $r$ iterations are used, the ideal good
probability is $\kappa=\sin^2((2r+1)\theta)$. A fully coherent
implementation can avoid postselection by defining
$W=\mathbf1_{\mathrm{good}}v$ and estimating
$\mathbb E W=\kappa\mathbb E_{q_s}v$, then dividing by a certified
positive $\kappa$. This includes failure branches in the observable,
charges the success predicate, and requires a rigorous error enclosure
for the moment, $\eta$, the chosen integer $r$ and $\kappa$. Treating
the prepared state as exactly $q_s$ without this correction is invalid.
Alternatively, generalised final amplification phases or conditional
postselection need their own gates, norm certificates and reset costs.

\begin{proposition}[Finite-precision cancellation cost]
Suppose an emitted $\widetilde B$ satisfies
$\|\widetilde B-B\|\leq\zeta<\sqrt\eta$. Conditional on a good
selector observation, its prepared state differs from the ideal
normalised good state by at most $2\zeta/\sqrt\eta$ in vector norm.
For any observable with operator norm at most one, the monetary
correction bias is at most
\[
 4b_sM_{2s}\frac\zeta{\sqrt\eta}
   =4b_sL\sqrt{M_{2s}}\zeta.
\]
For a fixed $r$-iteration coherent amplification tape, the operator-norm
error is bounded by $(2r+1)\zeta$ plus any reflection or phase errors.
\end{proposition}

\begin{proof}
Let $u$ and $\widetilde u$ be the two projected output vectors.
Projection is contractive, hence
$\|u-\widetilde u\|\leq\zeta$ and
$|\|u\|-\|\widetilde u\||\leq\zeta$.
Writing the difference of the two normalised vectors using the
denominator $\|u\|=\sqrt\eta$ gives a bound equal to the sum of these
two errors divided by $\sqrt\eta$. The conditional state exists since
$\|\widetilde u\|\geq\sqrt\eta-\zeta>0$.
The expectation difference for a norm-one observable between two unit
vectors is at most twice their vector distance, by insertion and
subtraction of one bra or ket at a time. Multiplication by the monetary
coefficient proves the bias bound. The amplification tape contains
$2r+1$ occurrences of $B$ or its adjoint; telescoping products of
unitaries sum their operator errors, proving the final assertion.
\end{proof}

\subsection{Weak-dispersion scaling and the required opponent}

Assume $d_i(\delta)=\delta a_i+O(\delta^2)$ and
$v_i(\delta)=\delta b_i+O(\delta^2)$, with the weighted centring
identities holding at every $\delta$. Suppose
$Q(y)=\tfrac12\sum_iw_i(a_i+b_i\cdot y)^2$ is not identically zero.
Expanding the exponentials with their integral remainders gives
$U(\delta,y)/\delta^2\to Q(y)$. On a fixed finite grid convergence
is uniform. On the Gaussian space a dominating envelope of the form
$C(1+\|y\|^2)e^{C\delta_0\|y\|}$ follows from
$|e^x-1-x|\leq x^2e^{|x|}/2$ and the centring identities.
Its every finite Gaussian moment is finite. Dominated convergence
therefore gives
\[
 M_{2s}=\Theta(\delta^{4s}),\quad
 L\longrightarrow2^s,\quad
 \eta=\Theta(\delta^{4s}).
\]
These statements concern the displayed labelled LCU; they are not
lower bounds for every possible representation of the importance state.

Repeated rejection preparation has expected cost $\Theta(\delta^{-4s})$
uses of $B$. Amplification can reduce the preparation-query order to
$\Theta(\delta^{-2s})$, provided all component preparation,
reflections and finite-precision errors are charged. Ideal amplitude
estimation of the normalised signed correction needs order
$b_sM_{2s}/\varepsilon$ calls at fixed confidence in the active regime.
The resulting preparation-query order is
$\delta^{2s}/\varepsilon$, before financial observable costs.
With rejection preparation the same calculation is order
$1/\varepsilon$. If $b_sM_{2s}\leq\varepsilon$, the deterministic
certificate already suffices and no correction estimator is necessary.

On the certified finite support an essential comparator is stronger.
Let $U_{\max}$ be a checked support bound and
$\rho_s=b_sU_{\max}^{2s}$. Estimating the direct signed observable
$R/\rho_s$ with the original base preparation needs order
$\rho_s/\varepsilon$ queries. For a fixed finite weak-dispersion
family $U_{\max}=O(\delta^2)$, hence this order is
$\delta^{4s}/\varepsilon$. It avoids the LCU preparation penalty.
Thus the proposed importance preparation is asymptotically worse in
this regime when oracle evaluation costs are comparable. A contrary
complete-cost conclusion requires a separately established saving in
the observable oracle, support certificate or another input regime.
The division by $U^{2s}$ in the importance observable can add substantial
fixed-point cost and conditioning error near zero; its definition at
zero must be implemented on all numerical inputs.

\subsection{No-arbitrage interface and the finite toy}

Let $\mathbb Q$ be the stipulated classical terminal measure and let
$D>0$ be the deterministic discount factor. I price each integrable
terminal cash flow $F$ by $\Pi(F)=D\mathbb E_{\mathbb Q}F$.
This is linear and positive, and $\Pi(1)=D$. If a traded forward
delivery is $\overline S_i$, I require
$\mathbb E_{\mathbb Q}S_i=\overline S_i$, so that the zero-cost
forward cash flow $S_i-\overline S_i$ has price zero. A nonnegative
terminal portfolio which is strictly positive on an event of positive
$\mathbb Q$ probability then has strictly positive price. Consequently
it cannot be acquired for nonpositive cost within this common pricing
functional. Quantum importance reweighting preserves this conclusion
only through the exact correction identity, not by treating $q_s$ as
a different financial pricing measure.

In the retained toy, each of two independent outer coordinates has
points $(-3/2,-1/2,1/2,3/2)$ with probability proportional to
$e^{-y^2/2}$. With independent $Z\sim N(0,1)$ and
$z_1(\delta)=\mathbb E e^{\delta Y_1}$, I set
\[
 S_\pm=K\exp\{\mu\pm\delta(Y_1+Y_2)+\lambda Z\},\qquad
 \mu=\log(\overline S/K)-\lambda^2/2-2\log z_1(\delta).
\]
Symmetry and independence give
$\mathbb E S_\pm=\overline S$ exactly. Their equal-weight basket
has ratio $T=\cosh(\delta(Y_1+Y_2))$ to its geometric basket.
The conditional capped payoff and its first Taylor coefficient are
computed in this same classical model; the retained numerical identity
checks the polynomial price plus signed correction against all sixteen
finite outer evaluations. It is a signal-construction test, not an
empirical sample or a replacement for the paper's continuous GBM model.

At $s=1$ and $\delta=1/4$, the exact finite LCU has three components,
coefficients $(-1,z_1(2\delta)/2,z_1(2\delta)/2)$, and success
probability approximately $0.00153885$. Its emitted signal circuit uses
eight qubits, sixty exact Toffolis and forty-two single-qubit rotation
uses. Eleven distinct retained rotation tapes have independent Acb
operator-norm certificates; the complete preparation has $4416$ T
gates and norm-error upper bound approximately $2.215\times10^{-8}$.
Nineteen ordinary amplification iterations use thirty-nine preparations
or adjoints. Charging the full zero reflection increases this to ten
qubits and $173288$ T gates, with an aggregate norm-error upper bound
approximately $8.637\times10^{-7}$. Literal gate replay gives good
probability approximately $0.99836037$ and clean work. Floating replay
is a diagnostic; the retained matrix enclosures supply gate-error rigour.

The remainder evaluator, its signed flag, deterministic representation
error, the end-to-end estimator, confidence postprocessing, readout and
reset are not yet supplied by this toy. I therefore make no complete
pricing-cost or advantage claim. Approximate prices must retain certified
intervals and the same risk-neutral expectation target. Independent
rounded point estimates can violate joint arbitrage constraints even
when each is accurate within its tolerance; the exact underlying
functional, or a jointly feasible price vector within the certified
intervals, is required for an across-contract no-arbitrage claim.
